The attention mechanism---the computational core of transformer architectures---is widely regarded as a breakthrough specific to natural language processing and large language models. We challenge this narrow view and investigate whether attention is a \emph{universal computational principle} that extends across all generative AI domains and, more broadly, underlies core internet infrastructure such as search ranking and recommendation systems. We conduct cross-domain experiments comparing attention-based models against non-attention baselines on text classification (\imdb) and image classification (\cifarten), analyze attention weight distributions across modalities, and formally demonstrate the mathematical isomorphism between transformer self-attention, PageRank, and collaborative filtering. Our results show that attention-based transformers outperform \bilstm baselines on text (73.7\% vs.\ 71.1\%, $p{=}0.046$), and that pretrained attention models dominate across modalities (\distilbert: 89.0\% on \imdb; \vit: 62.6\% vs.\ \resnet: 40.6\% on \cifarten). We further show that all three mechanisms---transformer attention, PageRank, and collaborative filtering---instantiate the same abstract operation: $\text{output} = \text{normalize}(\text{relevance}(\vq, \mK)) \cdot \mV$. These findings suggest that attention, understood as selective information routing through a capacity bottleneck, is a unifying principle across AI architectures and internet systems.
